% GENERATED FILE. Edit canonical lessons, then run npm run generate:books.
% canonical-source-hash: fnv1a64:df92e329f5b1860a
% canonical-lessons: AR-C58-qama, AR-C58-masha, AR-C58-rakada, AR-C58-laiba, AR-C58-ishtara

\chapter{To stand up, To walk, To run, To play, To buy}
\label{ch:ar-to-stand-up-to-walk-to-run-to-play-to-buy}
\hlchaptermodality{58}

\begin{chapteropening}
\textbf{By the end of this chapter:} \emph{I can say to stand up, to walk, to run, to play and to buy.}

Complete the last lesson of chapter 58: I can say to stand up, to walk, to run, to play and to buy.
\end{chapteropening}

\section[\texorpdfstring{\ar{قام} (qāma)}{قام (qāma)}]{\ar{قام} (qāma) — to stand up (he stood up)}

\label{lesson:AR-C58-qama}



\begin{quote}
Before the new one: say the Arabic for to sit, then the Arabic for to sleep.
\end{quote}

\subsection*{You'll want to know: \ar{قام}}
\textbf{\ar{قام}} — \emph{qāma} — \textquotedblleft{}to stand up (he stood up)\textquotedblright{}.

\begin{cousinweb}
\textbf{\ar{قام}} (\emph{qāma}) is built on the native Arabic root \emph{q-w-m}, \textquotedblleft{}to stand up, to rise\textquotedblright{} (Lane). The same root gives \emph{qawm}, \textquotedblleft{}a people\textquotedblright{}, and \emph{mustaqīm}, \textquotedblleft{}straight\textquotedblright{}.
\end{cousinweb}

\begin{grammarlens}[title={What you've built}]
One more word for this chapter.
\end{grammarlens}

\subsection*{Your turn}
\begin{itemize}
\raggedright
  \item \emph{Say it:} \emph{qāma}
  \item \emph{Say it:} \emph{qāma}, once more
  \item \emph{Say it:} say \emph{qāma}
  \item \emph{Recall:} say the Arabic for to sit, then the Arabic for to sleep, then say \emph{qāma} again
\end{itemize}

\begin{tcolorbox}[breakable,colback=teal!4,colframe=teal!35!black,title={Before you move on}]
What does \ar{قام} mean? (\textquotedblleft{}To stand up (he stood up)\textquotedblright{}.) Say it once more.
\end{tcolorbox}
\section[\texorpdfstring{\ar{مشى} (mashā)}{مشى (mashā)}]{\ar{مشى} (mashā) — to walk (he walked)}

\label{lesson:AR-C58-masha}



\begin{quote}
Before the new one: say the Arabic for to sleep, then the Arabic for to stand up.
\end{quote}

\subsection*{You'll want to know: \ar{مشى}}
\textbf{\ar{مشى}} — \emph{mashā} — \textquotedblleft{}to walk (he walked)\textquotedblright{}.

\begin{cousinweb}
\textbf{\ar{مشى}} (\emph{mashā}) is built on the native Arabic root \emph{m-sh-y}, \textquotedblleft{}to walk\textquotedblright{} (Lane). \emph{Mashā} is the plain verb, \textquotedblleft{}he walked\textquotedblright{}.
\end{cousinweb}

\begin{grammarlens}[title={What you've built}]
One more word for this chapter.
\end{grammarlens}

\subsection*{Your turn}
\begin{itemize}
\raggedright
  \item \emph{Say it:} \emph{mashā}
  \item \emph{Say it:} \emph{mashā}, once more
  \item \emph{Say it:} say \emph{mashā}
  \item \emph{Recall:} say the Arabic for to sleep, then the Arabic for to stand up, then say \emph{mashā} again
\end{itemize}

\begin{tcolorbox}[breakable,colback=teal!4,colframe=teal!35!black,title={Before you move on}]
What does \ar{مشى} mean? (\textquotedblleft{}To walk (he walked)\textquotedblright{}.) Say it once more.
\end{tcolorbox}
\section[\texorpdfstring{\ar{ركض} (rakaḍa)}{ركض (rakaḍa)}]{\ar{ركض} (rakaḍa) — to run (he ran)}

\label{lesson:AR-C58-rakada}



\begin{quote}
Before the new one: say the Arabic for to stand up, then the Arabic for to walk.
\end{quote}

\subsection*{You'll want to know: \ar{ركض}}
\textbf{\ar{ركض}} — \emph{rakaḍa} — \textquotedblleft{}to run (he ran)\textquotedblright{}.

\begin{cousinweb}
\textbf{\ar{ركض}} (\emph{rakaḍa}) is built on the native Arabic root \emph{r-k-ḍ}, \textquotedblleft{}to strike with the foot; to run\textquotedblright{} (Lane). Lane's first sense is striking with the foot, as in urging on a horse; \textquotedblleft{}to run\textquotedblright{} grew from it.
\end{cousinweb}

\begin{grammarlens}[title={What you've built}]
One more word for this chapter.
\end{grammarlens}

\subsection*{Your turn}
\begin{itemize}
\raggedright
  \item \emph{Say it:} \emph{rakaḍa}
  \item \emph{Say it:} \emph{rakaḍa}, once more
  \item \emph{Say it:} say \emph{rakaḍa}
  \item \emph{Recall:} say the Arabic for to stand up, then the Arabic for to walk, then say \emph{rakaḍa} again
\end{itemize}

\begin{tcolorbox}[breakable,colback=teal!4,colframe=teal!35!black,title={Before you move on}]
What does \ar{ركض} mean? (\textquotedblleft{}To run (he ran)\textquotedblright{}.) Say it once more.
\end{tcolorbox}
\section[\texorpdfstring{\ar{لعب} (laʿiba)}{لعب (laʿiba)}]{\ar{لعب} (laʿiba) — to play (he played)}

\label{lesson:AR-C58-laiba}



\begin{quote}
Before the new one: say the Arabic for to walk, then the Arabic for to run.
\end{quote}

\subsection*{You'll want to know: \ar{لعب}}
\textbf{\ar{لعب}} — \emph{laʿiba} — \textquotedblleft{}to play (he played)\textquotedblright{}.

\begin{cousinweb}
\textbf{\ar{لعب}} (\emph{laʿiba}) is built on the native Arabic root \emph{l-ʿ-b}, \textquotedblleft{}to play\textquotedblright{} (Lane). \emph{Laʿiba} is the plain verb, \textquotedblleft{}he played\textquotedblright{}.
\end{cousinweb}

\begin{grammarlens}[title={What you've built}]
One more word for this chapter.
\end{grammarlens}

\subsection*{Your turn}
\begin{itemize}
\raggedright
  \item \emph{Say it:} \emph{laʿiba}
  \item \emph{Say it:} \emph{laʿiba}, once more
  \item \emph{Say it:} say \emph{laʿiba}
  \item \emph{Recall:} say the Arabic for to walk, then the Arabic for to run, then say \emph{laʿiba} again
\end{itemize}

\begin{tcolorbox}[breakable,colback=teal!4,colframe=teal!35!black,title={Before you move on}]
What does \ar{لعب} mean? (\textquotedblleft{}To play (he played)\textquotedblright{}.) Say it once more.
\end{tcolorbox}
\section[\texorpdfstring{\ar{اشترى} (ishtarā)}{اشترى (ishtarā)}]{\ar{اشترى} (ishtarā) — to buy (he bought)}

\label{lesson:AR-C58-ishtara}



\begin{quote}
Before the new one: say the Arabic for to run, then the Arabic for to play.
\end{quote}

\subsection*{You'll want to know: \ar{اشترى}}
\textbf{\ar{اشترى}} — \emph{ishtarā} — \textquotedblleft{}to buy (he bought)\textquotedblright{}.

\begin{cousinweb}
\textbf{\ar{اشترى}} (\emph{ishtarā}) is built on the native Arabic root \emph{sh-r-y}, \textquotedblleft{}to sell; to buy\textquotedblright{} (Lane). Lane notes the base verb \emph{sharā} could mean both \textquotedblleft{}to buy\textquotedblright{} and \textquotedblleft{}to sell\textquotedblright{}; Form VIII \emph{ishtarā} settled on buying.
\end{cousinweb}

\begin{grammarlens}[title={What you've built}]
One more word for this chapter.
\end{grammarlens}

\subsection*{Your turn}
\begin{itemize}
\raggedright
  \item \emph{Say it:} \emph{ishtarā}
  \item \emph{Say it:} \emph{ishtarā}, once more
  \item \emph{Say it:} say \emph{ishtarā}
  \item \emph{Recall:} say the Arabic for to run, then the Arabic for to play, then say \emph{ishtarā} again
\end{itemize}

\begin{tcolorbox}[breakable,colback=teal!4,colframe=teal!35!black,title={Before you move on}]
What does \ar{اشترى} mean? (\textquotedblleft{}To buy (he bought)\textquotedblright{}.) Say it once more.
\end{tcolorbox}
